\subsection{Based saturation of a single weight in form degree zero}

\paragraph{A-S09-14 (Proved; the torsion-$H^1$ residue and its all-simplex repair).}
Fix a weight $w>0$ and identify its coefficient group
$S^w$ with $\mathbb Z$ by $T^{[w]}$.
Let $X$ be a simplicial nonnegative cochain complex of
free abelian groups with its constant coefficient map
$\underline{\mathbb Z}[0]\to X$.
Assume that its simplex-zero complex is exactly
$X_0=\mathbb Z[0]$, and that
$H^0(X_n)=\mathbb Z$ via the coefficient map for every
simplex $n$.  Let
\[
                 f:X\longrightarrow A,
                 \qquad g:X\longrightarrow B
 \tag{S09.36}
\]
be coefficient-preserving simplicial cochain maps.
In every simplex assume
$H^0(A_n)=H^0(B_n)=\mathbb Z$ via coefficients and
$H^1(A_n)=H^1(B_n)=0$.
The terms of $A$ and $B$ need not be flat.

There is an explicit simplicial cochain extension
$X\hookrightarrow\widetilde X$, with unchanged terms
in positive form degrees and free terms in degree zero,
to which both maps in (S09.36) extend.  It has
\[
 H^0(\widetilde X_n)=\mathbb Z,\qquad
 H^1(\widetilde X_n)=H^1(X_n)/\operatorname{tors}H^1(X_n),
 \qquad
 H^q(\widetilde X_n)=H^q(X_n)\quad(q\geq2).
 \tag{S09.37}
\]
This is a statement about the specified single-weight
linear complexes.  It neither kills free $H^1$ nor
asserts a full admissible PD-algebra completion.

\begin{proof}
For a vertex $v$ of $[n]$, write $a_v$ for the image
of $a\in X_n^0\otimes\mathbb Q$ in
$X_0^0\otimes\mathbb Q=\mathbb Q$.
Define
\[
 \widetilde X_n^0=
 \{a\in X_n^0\otimes\mathbb Q:
                  da\in X_n^1,\ a_0\in\mathbb Z\},
 \qquad
 \widetilde X_n^q=X_n^q\quad(q>0).
 \tag{S09.38}
\]
Since $X_n^2$ is torsion free, $c=da\in X_n^1$ is
an integral cycle.  Choose a primitive $h\in A_n^0$
of $f(c)$.  Every vertex of $c$ is zero because
$X_0^1=0$.  Consequently the vertex-zero value of
$h$ is a coefficient cycle, say $s\in\mathbb Z$.
The element
\[
                h_f(c)=h-s
 \tag{S09.39}
\]
is the unique primitive of $f(c)$ whose vertex-zero
value is zero.  Existence was just shown.  For
uniqueness, the difference of two primitives is a
coefficient cycle in $A_n^0$, and its zero vertex
value forces that coefficient to be zero.  Thus
$h_f$ is additive and integer-linear.  Define
\[
                 \widetilde f(a)=h_f(da)+a_0.
 \tag{S09.40}
\]
The identical construction defines $\widetilde g$.

We verify both the simplicial stability of (S09.38)
and the naturality of these maps.  Rationalizing the
target is harmless for this verification: flatness
of $\mathbb Q$ gives $H^0(A_n\otimes\mathbb Q)=\mathbb Q$
even when individual terms of $A_n$ have torsion.
The difference between $f_{\mathbb Q}(a)$ and
$h_f(da)$ is a rational coefficient cycle; evaluating
at vertex zero identifies it as $a_0$.  Therefore
\[
            f_{\mathbb Q}(a)=h_f(da)+a_0
                      \quad\text{in }A_n^0\otimes\mathbb Q.
 \tag{S09.41}
\]
At every other vertex $v$, the right side has an
integral coefficient value: its differential at that
vertex is zero and $H^0(A_0)=\mathbb Z$.
The left side has coefficient value $a_v$.
Since $\mathbb Z\to\mathbb Q$ is injective, $a_v$ is
an integer.  Thus the condition in (S09.38) is
independent of the chosen vertex.  Every simplicial
operator preserves it and preserves the condition
$da\in X_n^1$.  This makes $\widetilde X$ a
simplicial cochain complex containing $X$.

For a simplicial operator $\eta$, the two elements
$\eta\widetilde f(a)$ and $\widetilde f(\eta a)$
have the same differential.  Their difference is
therefore an integral coefficient cycle.  Equation
(S09.41) says their rationalizations agree; injectivity
on this coefficient group proves their equality
before rationalization.  The same argument applies
to $\widetilde g$.  Formula (S09.40) extends $f$ on
$X_n^0$, since for $a\in X_n^0$ the element
$f(a)-a_0$ is already its unique based primitive.
Thus the extensions are actual simplicial cochain
maps, including in targets with torsion.

The map
\[
       \widetilde X_n^0\longrightarrow
               \mathbb Z\oplus X_n^1,
                   \qquad a\longmapsto(a_0,da)
 \tag{S09.42}
\]
is injective.  Its kernel is a rational closed
zero-form, hence a rational coefficient, whose
vertex value is zero.  A subgroup of a free abelian
group is free, so (S09.42) proves the asserted
freeness, and finite freeness when $X_n^1$ is finite.
The same argument gives $H^0(\widetilde X_n)=\mathbb Z$.

Finally, $d\widetilde X_n^0$ is exactly the saturation
of $dX_n^0$ in $Z^1(X_n)$.  One inclusion follows
by clearing a denominator of $a$ in (S09.38).
For the other, let $c\in Z^1(X_n)$ and
$\ell c=db$ with $b\in X_n^0$ and $\ell\ne0$.
Then
\[
                 a=\frac{b-b_0}{\ell}
                       \in\widetilde X_n^0,
                 \qquad da=c.                       \tag{S09.43}
\]
This proves the assertion about $H^1$ in (S09.37).
All higher cohomology groups have unchanged defining
kernels and images, which proves the remaining claim.
\end{proof}

Here is the particular coefficient-residue calculation
inside this construction.  For $db=\ell c$ as in
(S09.43), any target primitive $h$ of the image of
$c$ has $\ell h-f(b)$ equal to a coefficient cycle.
Its coefficient modulo $\ell$ is exactly $-b_v$
at every chosen vertex $v$, for both target maps.
Indeed $h_v$ is itself an integral coefficient, so
evaluation gives $\ell h_v-b_v$.  After replacing
$h$ by its based primitive, the equality is exactly
$\ell h-f(b)=-b_0$.  Thus these two residues cannot
disagree when $X_0=\mathbb Z[0]$; no flatness or
unspecified comparison map is used in this argument.

For a positive-weight piece of a termwise torsion-free
strict PD algebra $D$ whose weight-zero part is exactly
$\mathbb Z$, the construction has a precise further use
when both targets are termwise torsion free.
Take the PD differential subalgebra of $D\otimes\mathbb Q$
generated by $D$ and the elements (S09.38) in that
weight.  The rationalizations of its two original
target maps send every new generator to the actual
integral elements (S09.40).  Since each target is
closed under its products, divided powers, and
differential, restriction gives two strict integral
PD maps from this subalgebra.  At the chosen weight,
products with a new positive-weight generator and
its higher divided powers occur in larger weights,
so its current-weight linear complex is exactly
the one just proved.  This verifies the stated
algebraic step of a weight induction.  No admissibility
or compatible later-weight repairs are inferred
from that step alone.
